\section{Introduction}\label{sec:intro}

A coin-turning walk on $\Z$ switches direction at step $k\ge2$ with
probability $c/k$, independently of the past, where $0<c<2$. This is the walk
of Engl\"ander and Volkov~\cite{ev2018,evw2020} with $p_k=c/k$. Let $A_N$ be
the number of right steps among the first $N$. With a fair first step
$A_N/N$ converges in law to $\mathrm{Beta}(c,c)$~\cite{ds2007,ev2018}. The
limit has a shape transition at $c=1$. Its density is U-shaped for $c<1$,
uniform at $c=1$ and peaked at $\frac12$ for $c>1$. We ask where this
transition happens at finite $N$, that is, at which $c$ an interior point of
$\{0,\dots,N\}$ becomes as likely as an end. With a fixed first step the
answer is exactly $c=1$, for every $N$ and every interior point
(Proposition~\ref{prop:aging-two}). With a fair first step the center ties
with the end at some $c_N<1$, and $1-c_N\sim\log2/\log N$
(Proposition~\ref{prop:aging-expansion}). The walk with $d$ states starts in
state~1, and at step $k\ge2$ it moves to each of the other $d-1$ states with
probability $\theta/(k+d-2)$. For $d=2$ this is $c/k$ with $c=\theta$. Its
count vector shows the same exact transition at $\theta=1$
(Theorem~\ref{thm:aging-d}).

\subsection{Setting and notation}\label{sec:notation}

The walk takes $N$ steps $X_1,\dots,X_N\in\{+1,-1\}$, and $X_N$ is the last
step. Let $A_N=\#\{i\le N:X_i=+1\}$ be the number of right steps, and write
$P_N(k)=\Prob(A_N=k)$ for the mass of the point $k$. We write $\Prob_+$ for
probability given $X_1=+1$. With a fair first step, $\Prob(X_1=+1)=\frac12$,
and $P_N(k)=P_N(N-k)$. In that case, for even $N$, we put $h=N/2$ and call $C_N=P_N(h)$
the center and $E_N=P_N(N)=P_N(0)$ the end. The points $1,\dots,N-1$ are interior. A crossing is a value of $c$ at
which an interior point ties with an end. Logarithms are natural, $H_k=\sum_{j\le k}1/j$,
and $\gamma$ is Euler's constant. Limits and $O(\cdot)$ refer to
$N\to\infty$, over even $N$ when $C_N$ appears.

\subsection{Results and method}\label{sec:results}

With a fixed first step and $c=1$, $A_N$ is exactly uniform on
$\{1,\dots,N\}$. The ratio of the mass of each interior point to that of the end increases
strictly in $c$. Hence every interior point ties with the end exactly at $c=1$, for every
$N\ge2$ (Proposition~\ref{prop:aging-two}). With $d$ states and $\theta=1$ the
count vector is exactly uniform over the compositions $(n_1,\dots,n_d)$ of $N$
with $n_1\ge1$, and every other count vector ties with the corner $(N,0,\dots,0)$
exactly at $\theta=1$ (Theorem~\ref{thm:aging-d}). Both laws come from a chain
that draws its state afresh or repeats it. Its counts and last state have an
explicit joint law for every weight (Proposition~\ref{prop:aging-urn}). With a
fair first step each interior point is twice as likely as an end at $c=1$, and
its crossing is unique and below~1 (Proposition~\ref{prop:aging-fair}). For
the center
$c_N=1-\log2/(\log N+\gamma+2-2\log2)+O((\log N)^{-3})$
(Proposition~\ref{prop:aging-expansion}, proved in Appendix~\ref{app:aging}).

The exact laws come from an induction that carries the last state along with
the counts. The ties come from writing each ratio of an interior point to the end as a sum
over paths of products of odds that increase in $c$. For the fair start we
couple the walk with the chain that draws afresh or repeats, whose counts are
a mixture of binomial laws. Summation by parts turns the coupling into a local
limit theorem with a rate at the center (Proposition~\ref{prop:D-llt}).

\paragraph{Verification.}\label{sec:verify}
No proof rests on a computation. Two scripts, \nolinkurl{verify_aging.py}
and \nolinkurl{verify_aging_urn.py}, recompute the numbers of this note,
and \nolinkurl{figures/make_aging_figure.py} draws Figure~\ref{fig:aging}. In exact
arithmetic $C_{200}-E_{200}$ changes sign between $c_{200}-10^{-6}$ and $c_{200}+10^{-6}$. Every
proposition, the lemma and the theorem are also proved in Lean~4 (version 4.33.1) with
Mathlib at commit \texttt{0df444a360ea}, in
\nolinkurl{kagey-problems/problem131/lean/PaperD} of
\url{https://github.com/apemm/Kagey-Problems}, using only the axioms
\texttt{propext}, \texttt{Classical.choice} and \texttt{Quot.sound}. The file
\nolinkurl{MAP.md} there lists them under the numbers of an earlier draft.
Its numbers 3.1 to 3.7 are 2.1 to 2.7 here, and its numbers B.1 to B.3 are
A.1 to A.3 here.

\subsection{Prior work}\label{sec:prior}

Dietz and Sethuraman~\cite{ds2007} proved the limit law of $A_N/N$,
Engl\"ander and Volkov~\cite{ev2018} gave a second proof, and Engl\"ander,
Volkov and Wang~\cite{evw2020} found the scaling limit. For $d$ states the
limit is Dirichlet~\cite[Theorem~1.4]{ds2007}, and Bouguet and
Cloez~\cite[Propositions~4.1 and~4.5]{bc2018} find the joint limit with the
current state. The uniform law for 2 states is an exercise
in~\cite[Section~3.1, footnote~1]{evw2020}. The count law of the chain of
Proposition~\ref{prop:aging-urn} also follows from Dietz, Lippitt and
Sethuraman~\cite[Theorem~2.18]{dls2023}. Chapter~7 of the March 18, 2024
author draft~\cite{ev2024} of Engl\"ander and Volkov's
monograph~\cite{ev2025} gives the two-state joint law with the last state for
integer urn weights. We cite the draft for every locator in that work.
Section~\ref{sec:aging} gives the precise overlaps. Pradhan, Roy and
Ray~\cite{prr2026} find, with the methods of physics, the second moment and
the run tail of the walk with switch probability $c/k^\alpha$, which is ours
at $\alpha=1$. What we add is the joint law for $d\ge3$ or a
weight that is not an integer, the strict ordering that turns the uniform laws
into ties (Proposition~\ref{prop:aging-two}(b) and
Theorem~\ref{thm:aging-d}(b)), and the expansion of $c_N$ with the local limit
theorem behind it. For runs of geometric length the center overtakes the
ends when the expected number of switches is about $\log N$, for 2
directions~\cite{paperA} and for $d$ directions~\cite{paperB}.
